\section{Discussion and Limitations}

The SECDED experiment shows why code semantics, architecture, and physical outcome must be separately named. Exact equivalence permits causal attribution to the evaluated temporal microarchitecture as a whole, while matched seeds test whether its trade direction is an artifact of one heuristic trajectory. The envelope is intentionally descriptive: five fixed seeds cannot estimate a universal distribution, tail probability, confidence interval, or tool-independent effect.

Hsiao recovery supplies a second lesson. The old decoder's policy failure concerned representation, not code quality. Replacing its inferred table with explicit comparisons preserved behavior and passed the unchanged flow, but created a new hardware identity whose PPA cannot be assigned to the table implementation. Conversely, exact equivalence is what permits the prior guarantee semantics to transfer without repeating an ambiguous code-level experiment.

Correction tiers must remain separate. Within the stated physical objectives, an SECDED point may be Pareto-superior to the evaluated BCH realization, yet it cannot correct every double-coordinate error. A system with mandatory W2 correction may accept a timing gap or seek another BCH/DEC architecture. The data therefore support conditional selection rather than a family-wide dominance statement. Weight-3 counts are omitted from the main results because their unweighted canonical masks have no demonstrated operational probability; they remain useful artifact-level decoder checks.

The quantitative scope is one open 130-nm library, one TT 1.80~V/25$^{\circ}$C corner, one 10~ns target, one floorplan policy, and specific RTL. OpenSTA activity power is an estimate rather than silicon measurement. We exclude SRAM macro/storage organization, PVT variation, IR drop, clock-domain integration, scrubbing, interleaving, physical bit adjacency, FIT/SER, and field-weighted SDC/DUE. Common policy reduces configuration confounding but does not prove optimal floorplanning for each architecture. These limits bound numerical generalization while leaving the within-flow identity and sensitivity results independently auditable.

\subsection{Selection Implications}

The evidence suggests a staged selection rule. First filter implementations by categorical guarantee, target feasibility, maximum request latency, and evidence status. Only then compare surviving hardware identities on area, routing, and achievable energy. A one-cycle latency requirement favors combinational SECDED over its exact-equivalent pipeline even though the latter has higher frequency and lower streaming energy. A strict W2-correction requirement admits the BCH semantics but exposes a physical implementation gap at 10~ns. Hsiao is now a measured same-guarantee alternative rather than an unavailable label, yet its aggregate outcome cannot be assigned to all Hsiao architectures.

This ordering also clarifies what future work could change. A multi-corner sweep could test whether timing and power orderings survive PVT variation; another library or node could change XOR/register cost ratios; a second qualified BCH decoder could show whether the present target miss is architectural; and a physical memory-fault model could weight beyond-guarantee masks. Each extension should create a new named physical or hardware identity rather than overwrite this dataset. None would alter the exact equivalence proof or the measurements of the frozen implementations.

\subsection{Claim Boundaries and Portability}

The evidence separates conclusions with different portability. Exact RTL equivalence is independent of the route seeds and establishes behavior for the frozen SECDED pair and Hsiao decoder relation. Exhaustive guarantee tests establish categorical behavior only for the qualified code/decoder universe. Timing, area, routing, and activity-power results are properties of the frozen hardware and physical identities; their numerical values should not be transferred to a different decoder, constraint set, library, corner, or tool revision. Finally, seed-ordering statements are properties of this closed five-seed sensitivity set. Making these levels explicit avoids both underclaiming the proofs and overclaiming the physical data.

Several tempting inferences are therefore excluded. The frequency metric does not say the designs were reimplemented at their maximum clock. Equal-work no-error energy does not predict energy under correction-heavy traffic. A DRC-clean BCH route does not imply target feasibility, and a target miss does not invalidate the BCH correction guarantee. Hsiao's successful algorithmic implementation does not establish that all Hsiao matrices or syndrome organizations have similar PPA. Likewise, stable ordering over five matched seeds is evidence against a one-seed artifact under the pinned heuristic controls, not evidence that another flow, node, or seed range must preserve it.

These boundaries still allow useful engineering portability. The identity/evidence method, missing-data rules, matched-pair construction, and categorical treatment of correction strength can be reused without assuming the present numerical effects. A replication can substitute another library, corner, or qualified implementation while retaining the same joins and decision logic, then compare whether signs, feasibility, and Pareto relations change. Reporting a divergence would be informative rather than contradictory because it would attach the new outcome to a newly named physical identity.

\subsection{Auditability}

The artifact retains full-precision records even though the paper rounds for display. Each route has a unique run ID, complete stage status, raw ORFS JSON and reports, metadata, and SHA-256 inventory. Post-execution analysis verifies every raw manifest before generating the per-seed CSV, architecture summaries, paired-effect JSON, and manuscript registry. Figures and tables are regenerated from that registry; an independent validator recomputes headline macros, rejects ambiguous timing terminology, and rechecks frozen Hsiao and Revision-1 hashes. This discipline makes each cross-layer claim independently auditable without making provenance itself the scientific contribution.
